\documentclass[a4paper,11pt]{article}
\usepackage[utf8]{inputenc}
\usepackage[german]{babel}
\usepackage{geometry}
\usepackage{array}
\usepackage{colortbl}
\usepackage{xcolor}
\usepackage{pdflscape}
\usepackage{longtable}

\geometry{a4paper, margin=1.5cm, landscape}

% Farben definieren
\definecolor{headercolor}{RGB}{69, 90, 100}

\begin{document}

\section*{Manuelle Tests}

\renewcommand{\arraystretch}{1.8}
\footnotesize
\begin{longtable}{|p{3.3cm}|p{3.3cm}|p{3.3cm}|p{5.5cm}|p{3.3cm}|}
	\hline
	\rowcolor{headercolor}
	\textcolor{white}{\textbf{Test ID}}                                                                                                                                & \textcolor{white}{\textbf{Ziel}} & \textcolor{white}{\textbf{Rolle}} & \textcolor{white}{\textbf{Beschreibung}} & \textcolor{white}{\textbf{Ergebnis}} \\
	\hline
	\endfirsthead

	\multicolumn{5}{c}%
	{{\tablename\ \thetable{} -- Fortsetzung}}                                                                                                                                                                                                                                                                                  \\
	\hline
	\rowcolor{headercolor}
	\textcolor{white}{\textbf{Test ID}}                                                                                                                                & \textcolor{white}{\textbf{Ziel}} & \textcolor{white}{\textbf{Rolle}} & \textcolor{white}{\textbf{Beschreibung}} & \textcolor{white}{\textbf{Ergebnis}} \\
	\hline
	\endhead

	\hline
	\multicolumn{5}{r}{{Fortsetzung auf nächster Seite}}                                                                                                                                                                                                                                                                        \\
	\endfoot

	\hline
	\caption{Tabelle der manuell durchgeführten Tests}
	\endlastfoot

	\textbf{T2.1.1}                                                                                                                                                    &
	Testen der User Story S2.1                                                                                                                                         &
	Nutzer                                                                                                                                                             &
	Als angemeldeter Nutzer gelange ich mit einem Klick auf den Button \textit{Create New User Story} auf die Seite, auf welcher man User Stories anlegen kann.        &
	Nicht bestanden
	\\
	\hline

	\textbf{T2.1.2}                                                                                                                                                    &
	Testen der User Story S2.1                                                                                                                                         &
	Nutzer                                                                                                                                                             &
	Als Nutzer kann ich auf der Seite \textit{/create-user-story} eine neue User Story anlegen.                                                                        &
	Nicht bestanden
	\\
	\hline

	\textbf{T4.1.1}                                                                                                                                                    &
	Testen der User Story S4.1                                                                                                                                         &
	Nutzer                                                                                                                                                             &
	Als angemeldeter Nutzer kann ich auf den Button \textit{Create New Backlog} klicken. Dann öffnet sich die Seite, auf welcher man ein neues Backlog erstellen kann. &
	Nicht bestanden
	\\
	\hline

	\textbf{T4.1.2}                                                                                                                                                    &
	Testen der User Story S4.1                                                                                                                                         &
	Nutzer                                                                                                                                                             &
	Als angemeldeter Nutzer kann ich auf der Seite \textit{/create-backlog} ein neues Backlog anlegen.                                                                 &
	Nicht bestanden
	\\
	\hline
\end{longtable}

\end{document}
